\documentclass[12pt,a4paper]{article}

\usepackage[T1]{fontenc}
\usepackage[utf8]{inputenc}
\usepackage{amsmath,amssymb,amsfonts}
\usepackage{graphicx}
\usepackage{hyperref}
\usepackage{xcolor}
\usepackage{booktabs}
\usepackage{enumitem}
\usepackage{float}
\usepackage{geometry}

\geometry{
  left=3.2cm,
  right=3.2cm,
  top=2.8cm,
  bottom=2.8cm
}

\hypersetup{
  colorlinks=true,
  linkcolor=blue!60!black,
  citecolor=blue!60!black,
  urlcolor=blue!60!black
}

\setlength{\parindent}{0pt}
\setlength{\parskip}{0.6em}

\pagestyle{plain}

\begin{document}

{\small\scshape paper}

\vspace{0.8cm}

{\LARGE The Platform Has a Food Chain. You Are In It.\\[6pt]
\normalsize An ecological model of attention with five rational agents}

\vspace{0.7cm}

Hakvin Vosteen

\vspace{0.6cm}

\noindent\rule{\textwidth}{0.4pt}

\textbf{Abstract}

\vspace{0.2em}

Every short-video platform produces an attention pool $A_0$ generated by creators and consumed by a small number of recurring downstream roles. We classify these roles biologically: Whale (creator, source), Remora (sharer, mutualist), Barnacle (hater, commensal), Plankton (lurker, neutral), Parasite (drama aggregator). Each agent extracts a multiplicative share $\alpha_i$ of the upstream attention. The total attention generated by a cascade of $N$ agents is $\mathcal{A}(A_0) = A_0 \sum_{i=0}^{N} \prod_{j=1}^{i} \alpha_j(s_j^{*})$, where each $s_j^{*}$ is the strategy chosen by agent $j$ to maximise its own extraction. We show that the platform-level objective of maximising $\mathcal{A}$ is equivalent to maximising the cascade depth, which under the empirical inequality $\alpha_{\mathrm{hate}} > \alpha_{\mathrm{agree}}$ produces conflict-amplifying recommendation as a Nash equilibrium of the cascade. The model gives a clean Lotka-Volterra coexistence between creators and haters, an interior optimum (not maximum) in conflict intensity for the hater agent, and a structural identification of which agents are extractive (parasite) versus load-bearing-but-invisible (plankton). Nobody in the system is irrational. The platform built the food chain.

\vspace{0.4em}

{\small\textbf{Keywords:} attention economics, recommender systems, ecological models, Lotka-Volterra dynamics, Nash equilibrium, platform design, social media}

\noindent\rule{\textwidth}{0.4pt}

\vspace{0.5cm}

% ============================================================
\section{Introduction}
\label{sec:intro}
% ============================================================

Every person you have ever seen on a social-media platform is one of five biological species. Five agents, one attention pool. The species are not metaphorical. They follow the same extraction rules, the same coexistence dynamics, and the same Nash equilibria as their ecological counterparts. The only thing different is the substrate: instead of carbon and nitrogen, the resource is attention $A_0$.

Section~\ref{sec:agents} introduces the five agents and their extraction coefficients. Section~\ref{sec:dynamics} solves the local dynamics for each agent in isolation. Section~\ref{sec:cascade} composes the agents into a cascade and writes the master equation. Section~\ref{sec:platform} shows that platform-level optimisation produces conflict amplification as a Nash equilibrium. Section~\ref{sec:predictions} states testable predictions. Section~\ref{sec:boundaries} discusses limitations.

% ============================================================
\section{Five agents on one attention pool}
\label{sec:agents}
% ============================================================

Let $A_0$ denote the attention generated by the creator at the head of a content cascade. Every downstream agent extracts a multiplicative share
\begin{equation}
  A_i = \alpha_i \cdot A_{i-1},
  \qquad
  \alpha_i \in (0, 1] \text{ for non-parasitic agents,} \quad \alpha_i > 1 \text{ for parasites.}
  \label{eq:extract}
\end{equation}
Nobody creates from zero. Every agent depends on the upstream node.

\begin{figure}[H]
  \centering
  \includegraphics[width=0.85\textwidth]{figures/fig1_ecology.pdf}
  \caption{The five-agent ecology. The creator is the source of $A_0$. Sharers (remoras) amplify it. Haters (barnacles) attach passively. Lurkers (plankton) drift in the wake. Parasites (drama aggregators) bore into the original output and extract $A_0 \cdot \alpha_D$ with $\alpha_D > 1$.}
  \label{fig:ecology}
\end{figure}

The five agents and their biological analogues are summarised in Table~\ref{tab:agents}.

\begin{table}[H]
\centering
\begin{tabular}{lllcc}
\toprule
Role             & Biology   & $\alpha_i$       & $\partial \mathrm{Harm} / \partial \alpha_i$ & Type     \\ \midrule
Creator (Whale)  & source    & $1.0$            & source                                       & source   \\
Sharer (Remora)  & mutualist & $0.05$ to $0.25$ & $< 0$                                         & mutualist \\
Hater (Barnacle) & commensal & $0.10$ to $0.80$ & $\approx 0$                                  & commensal \\
Lurker (Plankton)& passive   & $\approx 0$      & $0$                                          & neutral   \\
Drama (Parasite) & parasite  & $> 1$            & $> 0$                                        & parasite  \\ \bottomrule
\end{tabular}
\caption{The five-agent system. $\alpha_i$ is the extraction coefficient. The harm derivative measures whether increasing the agent's intensity harms the creator. Only the parasite has a positive sign.}
\label{tab:agents}
\end{table}

% ============================================================
\section{Per-agent dynamics}
\label{sec:dynamics}
% ============================================================

\subsection{Whale (Creator). Source. \texorpdfstring{$\alpha_W = 1.0$}{alpha\_W = 1.0}}

The creator generates $A_0$. Every downstream agent lives off this number. The creator is vulnerable to barnacle overload: once the population of haters $n_B$ exceeds a critical threshold $n_B^{*}$, the drag they impose, $D(n_B) = \alpha_B \cdot n_B$, exceeds the creator's output and the creator slows. The strategy set is $\mathcal{S}_W = \{\text{post}, \text{pause}, \text{respond}, \text{ignore}\}$.

\begin{figure}[H]
  \centering
  \includegraphics[width=0.78\textwidth]{figures/fig2_whale_growth.pdf}
  \caption{Creator output $A_0(t)$ over time. After an initial growth phase the system reaches a stable plateau. When the barnacle population crosses the critical threshold $n_B^{*}$ the drag exceeds output and the creator enters burnout (dashed segment).}
  \label{fig:whale_growth}
\end{figure}

The critical barnacle load is given by
\begin{equation}
  n_B^{*} = \frac{A_0}{\alpha_B},
  \label{eq:nstar}
\end{equation}
the maximum number of haters the creator can absorb before output collapses.

\subsection{Remora (Sharer). Mutualist. \texorpdfstring{$\alpha_R \in [0.05, 0.25]$}{alpha\_R in [0.05, 0.25]}}

The sharer swims with the whale. Amplifies $A_0$, removes spam, helps the creator. Timing matters. Share too late and the attention peak is gone. The strategy set is $\mathcal{S}_R = \{\text{early}, \text{peak}, \text{late}, \text{add context}\}$. The amplification function is
\begin{equation}
  A_R(\tau) = \alpha_R \, A_0 \, e^{-\lambda (\tau - \tau^{*})^2},
  \label{eq:remora}
\end{equation}
peaking at $\tau^{*} \approx 3$ hours after the original post.

\begin{figure}[H]
  \centering
  \includegraphics[width=0.78\textwidth]{figures/fig3_remora_timing.pdf}
  \caption{Remora amplification as a function of share time. The peak at $\tau^{*} \approx 3$ hours is roughly the half-life of attention on the typical platform; sharing earlier or later forfeits a multiplicative factor of attention.}
  \label{fig:remora}
\end{figure}

\subsection{Barnacle (Hater). Commensal. \texorpdfstring{$\alpha_B(c) = \alpha_0 + \beta c^{\gamma}, \gamma < 1$}{alpha\_B(c) = alpha\_0 + beta c\^gamma, gamma < 1}}

The hater attaches and filters passively. Not evil. Not broken. Rational. The empirical regularity $\alpha_{\mathrm{hate}} > \alpha_{\mathrm{agree}}$ holds because conflict generates more secondary engagement loops than agreement does. Hating is therefore the Nash equilibrium strategy in the secondary-engagement subgame. The strategy set is $\mathcal{S}_B = \{\text{agree}, \text{critique}, \text{ratio}, \text{insult}, \text{strawman}\}$.

The extraction coefficient as a function of conflict intensity $c$ has the form
\begin{equation}
  \alpha_B(c) = \alpha_0 + \beta c^{\gamma}, \qquad \gamma < 1,
  \label{eq:barnacle}
\end{equation}
which is concave in $c$. The interior optimum is
\begin{equation}
  c^{*} = \left( \frac{\alpha_0}{\beta \gamma} \right)^{1/(\gamma - 1)},
  \label{eq:cstar}
\end{equation}
which is strictly less than the maximum conflict intensity. The platform-optimal hater therefore tunes conflict to an interior level rather than maximum, contradicting the naive view that engagement-maximising platforms push conflict to the extreme.

\begin{figure}[H]
  \centering
  \includegraphics[width=0.78\textwidth]{figures/fig4_barnacle_conflict.pdf}
  \caption{Hater extraction $\alpha_B(c)$ as a function of conflict intensity. Concave with $\gamma < 1$. The interior optimum $c^{*}$ is strictly below maximum conflict; the equilibrium is in the ratio regime, not the strawman regime.}
  \label{fig:barnacle}
\end{figure}

\subsection{Plankton (Lurker). Passive. \texorpdfstring{$\alpha_P \approx 0$}{alpha\_P approx 0}}

The lurker drifts in the whale's wake. Consumes without reacting. Invisible to the platform's scoring function. Invisible to the creator. Yet essential: without lurkers there is no audience, no creator, no barnacles. The strategy set is $\mathcal{S}_P = \{\text{watch}, \text{scroll}, \text{re-watch}\}$.

The harm derivative vanishes:
\begin{equation}
  \alpha_P \approx 0, \qquad \frac{\partial \mathrm{Harm}}{\partial \alpha_P} = 0,
  \label{eq:plankton}
\end{equation}
making the plankton invisible but load-bearing.

\subsection{Parasite (Drama Aggregator). \texorpdfstring{$\alpha_D > 1$}{alpha\_D > 1}}

The parasite is the only agent in the ecosystem that extracts more than the original post generated. It bores into the whale, builds threads and articles and videos and podcasts about the conflict, and the whale loses narrative control permanently. The harm derivative is positive: $\partial \mathrm{Harm} / \partial \alpha_D > 0$. The strategy set is $\mathcal{S}_D = \{\text{thread}, \text{article}, \text{video}, \text{podcast}\}$. The parasite is structurally distinct from the other four agents because $\alpha_D > 1$, violating the local conservation that holds for the other agents.

% ============================================================
\section{The cascade and the master equation}
\label{sec:cascade}
% ============================================================

The total attention generated by a cascade of $N$ agents acting on the seed $A_0$ is
\begin{equation}
  \boxed{\quad \mathcal{A}(A_0) \;=\; A_0 \sum_{i=0}^{N} \prod_{j=1}^{i} \alpha_j(s_j^{*}) \quad}
  \label{eq:master}
\end{equation}
where $s_j^{*}$ is the strategy chosen by agent $j$ to maximise its own extraction:
\begin{equation}
  s_j^{*} = \arg \max_{s \in \mathcal{S}_j} \alpha_j(s) \cdot A_{j-1}, \qquad \forall j.
  \label{eq:nash}
\end{equation}
The simultaneous solution of \eqref{eq:nash} across all agents is the Nash equilibrium of the cascade.

The harm-weighted susceptibility of agent $i$ to the upstream creator is
\begin{equation}
  \kappa_i = \frac{\partial A_i / \partial A_0}{\partial \mathrm{Harm} / \partial A_0},
  \label{eq:kappa}
\end{equation}
which gives a normalised measure of how much each agent's extraction scales with creator output relative to the harm done. Mutualists have $\kappa < 0$ (they benefit the creator), commensals have $\kappa \approx 0$, and parasites have $\kappa > 0$ unbounded.

% ============================================================
\section{Lotka-Volterra coexistence}
\label{sec:lv}
% ============================================================

Coupling the creator population $W(t)$ with the hater population $B(t)$ produces a standard Lotka-Volterra system:
\begin{align}
  \dot{W} &= \alpha W - \beta W B, \label{eq:lv1} \\
  \dot{B} &= \delta W B - \gamma B. \label{eq:lv2}
\end{align}
The interior fixed point is
\begin{equation}
  W^{*} = \frac{\gamma}{\delta}, \qquad B^{*} = \frac{\alpha}{\beta},
  \label{eq:lv_eq}
\end{equation}
giving stable coexistence as a digital ecology.

\begin{figure}[H]
  \centering
  \includegraphics[width=0.85\textwidth]{figures/fig5_lotka_volterra.pdf}
  \caption{Numerical solution of the creator-hater Lotka-Volterra system \eqref{eq:lv1}-\eqref{eq:lv2}. Standard predator-prey oscillation with hater density lagging creator density. Neither population goes to zero. The platform-level cascade is built on top of this stable coexistence.}
  \label{fig:lv}
\end{figure}

The Lotka-Volterra dynamics are local. The full cascade adds the remora, plankton, and parasite agents on top, but the creator-hater pair is the load-bearing oscillator of the system; the other three agents can be analysed in quasi-steady-state relative to the $W$-$B$ cycle.

% ============================================================
\section{Why the platform pushes conflict}
\label{sec:platform}
% ============================================================

The platform maximises $\mathcal{A}(A_0)$ over all cascade configurations. From the master equation \eqref{eq:master}, this is equivalent to maximising the product $\prod_j \alpha_j(s_j^{*})$, which under the inequality $\alpha_{\mathrm{hate}} > \alpha_{\mathrm{agree}}$ implies that the equilibrium strategy at the hater node is conflict, not agreement. The platform pushes conflict not because it wants harm but because conflict is the local maximum of the cascade.

\textbf{Key insight.} Nobody in this ecosystem is irrational. Every agent is executing the optimal strategy given the incentive structure the platform designed. The platform built the food chain. You are somewhere in it.

This explains a phenomenon that the standard "users are addicted" narrative cannot: the same user behaves rationally on a platform with $\alpha_{\mathrm{hate}} > \alpha_{\mathrm{agree}}$ and rationally on a platform without that inequality, but the equilibrium is qualitatively different in the two cases. The behaviour is downstream of the incentive structure, not upstream of it.

% ============================================================
\section{Testable predictions}
\label{sec:predictions}
% ============================================================

\textbf{Prediction 1} (Interior conflict optimum). Empirically estimated $\alpha_B(c)$ on a representative platform sample should be concave in $c$ with an interior maximum strictly below the upper bound of the conflict scale. Falsified if the data show a monotone-increasing or monotone-decreasing $\alpha_B(c)$.

\textbf{Prediction 2} (Mutualist timing). The optimal share-time $\tau^{*}$ for sharers should equal the half-life of the upstream content's attention curve. Falsified if remoras succeed equally well at all times after posting.

\textbf{Prediction 3} (Parasite distinguishability). Drama aggregators should show $\alpha_D > 1$ in revealed-engagement metrics, distinguishable from haters by a strictly positive secondary-content-creation rate. Falsified if the two roles cannot be separated empirically.

\textbf{Prediction 4} (Burnout threshold). Creator burnout should track \eqref{eq:nstar}: the burnout time should be inversely proportional to the empirical hater extraction coefficient $\alpha_B$ controlling for $A_0$. Falsified if burnout is independent of the hater population given $A_0$.

% ============================================================
\section{Boundaries}
\label{sec:boundaries}
% ============================================================

The five-agent classification is empirically derived but not unique. A finer taxonomy might distinguish sub-types within the remora (e.g.\ stitcher versus duet partner) or within the parasite (e.g.\ aggregator versus reaction-creator). The cascade depth $N$ is treated as fixed; in reality it is bounded by the platform's recommendation graph and may be endogenous.

The harm function is left implicit. A full treatment requires specifying $\mathrm{Harm}(\alpha_i, A_0)$ as a function of agent intensity and creator output, which is platform-dependent and time-varying.

The Lotka-Volterra coupling assumes that haters depend only on creators and creators depend only on haters. In reality the parasite agent couples back into the dynamics with $\alpha_D > 1$ and may destabilise the otherwise stable coexistence. A three-species generalisation is the natural next step.

% ============================================================
\section{Conclusion}
\label{sec:conclusion}
% ============================================================

The platform has a food chain. It has five species. Each species has a strategy set. Each strategy is rational given the incentive structure. The platform-level objective maximises the product of extraction coefficients along the cascade, and under the empirical inequality $\alpha_{\mathrm{hate}} > \alpha_{\mathrm{agree}}$, the equilibrium pushes conflict to an interior optimum, not zero and not maximum. Nobody is irrational. The platform built the food chain. You are somewhere in it.

\vspace{1.2em}

\textbf{Code.} Replication code at \url{https://github.com/hakvinv/platform-ecology}.

% ============================================================
\begin{thebibliography}{99}

\bibitem{simon1971designing}
H.~A. Simon,
\textit{Designing organizations for an information-rich world},
in M.~Greenberger (ed.), Computers, Communications, and the Public Interest, Johns Hopkins University Press (1971) 37--72.

\bibitem{lotka1925elements}
A.~J. Lotka,
\textit{Elements of Physical Biology},
Williams \& Wilkins, Baltimore (1925).

\bibitem{volterra1926fluctuations}
V.~Volterra,
\textit{Fluctuations in the abundance of a species considered mathematically},
Nature \textbf{118} (1926) 558--560.

\bibitem{wu2007novelty}
F.~Wu, B.~A. Huberman,
\textit{Novelty and collective attention},
PNAS \textbf{104} (2007) 17599--17601.

\bibitem{vosipoulos2014echo}
E.~Bakshy, S.~Messing, L.~A. Adamic,
\textit{Exposure to ideologically diverse news and opinion on Facebook},
Science \textbf{348} (2015) 1130--1132.

\bibitem{munger2020rightwing}
K.~Munger, J.~Phillips,
\textit{Right-wing YouTube: a supply and demand perspective},
The International Journal of Press/Politics \textbf{27} (2020) 186--219.

\bibitem{vosteen2026reels}
H.~Vosteen,
\textit{Instagram Reels Is a TikTok Echo: a Stackelberg Imitation Model of Platform Convergence},
Working paper (2026).

\end{thebibliography}

\end{document}
